\documentclass[12pt,a4paper]{article}

% =========================================================
% PAQUETES
% =========================================================
\usepackage[utf8]{inputenc}
\usepackage[T1]{fontenc}
\usepackage[spanish,es-nodecimaldot]{babel}
\usepackage{lmodern}
\usepackage{geometry}
\usepackage{graphicx}
\usepackage{booktabs}
\usepackage{array}
\usepackage{longtable}
\usepackage{float}
\usepackage{amsmath,amssymb}
\usepackage{enumitem}
\usepackage{xcolor}
\usepackage{titlesec}
\usepackage{fancyhdr}
\usepackage{hyperref}
\usepackage{setspace}
\usepackage{lastpage}
\usepackage{caption}
\usepackage{subcaption}
\usepackage{tikz}
\usepackage{tabularx}
\usepackage{listings}

% =========================================================
% CONFIGURACIÓN GENERAL
% =========================================================
\geometry{
    top=2.5cm,
    bottom=2.5cm,
    left=2.7cm,
    right=2.7cm
}

\setstretch{1.15}
\setlength{\parindent}{0pt}
\setlength{\parskip}{0.6em}

% Color institucional aproximado para títulos.
\definecolor{UBUblue}{RGB}{0,82,147}
\definecolor{lightgray}{RGB}{245,245,245}
\definecolor{UBUlight}{RGB}{235,243,249}
\definecolor{UBUdark}{RGB}{34,45,58}

% Carpeta donde se encuentra el logotipo institucional
\graphicspath{{logo/}}

\hypersetup{
    colorlinks=true,
    linkcolor=UBUblue,
    urlcolor=UBUblue,
    citecolor=UBUblue,
    pdftitle={Practicas de laboratorio - Hardware de Aplicacion Especifica},
    pdfauthor={Universidad de Burgos}
}

% =========================================================
% DATOS DEL TRABAJO
% =========================================================

\newcommand{\universidad}{Universidad de Burgos}
\newcommand{\centro}{Escuela Politécnica Superior}
\newcommand{\grado}{Grado en Ingeniería Informática}
\newcommand{\asignatura}{Hardware de Aplicación Específica}
\newcommand{\titulotrabajo}{Prácticas de laboratorio}
\newcommand{\subtitulotrabajo}{Bloque 1}

\newcommand{\estudianteA}{Nombre y apellidos 1}
\newcommand{\estudianteB}{Nombre y apellidos 2}
\newcommand{\correo}{correo@alu.ubu.es}
\newcommand{\grupo}{Grupo}
\newcommand{\profesor}{César Represa Pérez}
\newcommand{\cursoacademico}{2026--2027}
\newcommand{\fechaentrega}{\today}

% =========================================================
% COMANDOS AUXILIARES
% =========================================================

% Figura con captura de MATLAB. Si el fichero no existe se
% dibuja un recuadro indicando el nombre que debe tener.
\newcommand{\captura}[3]{%
    \begin{figure}[H]
        \centering
        \IfFileExists{#1}{%
            \includegraphics[width=0.92\textwidth]{#1}%
        }{%
            \fbox{%
                \parbox[c][6.5cm][c]{0.86\textwidth}{%
                    \centering
                    Captura pendiente de insertar\\[0.3cm]
                    \texttt{#1}%
                }%
            }%
        }
        \caption{#2}
        \label{#3}
    \end{figure}%
}

% Bloque para la respuesta a una cuestión del enunciado.
\newenvironment{respuesta}
    {\par\medskip\noindent
     \begin{minipage}{\linewidth}
     \small
     \hrule\vspace{0.4em}\noindent\textbf{Respuesta.}\ }
    {\vspace{0.4em}\hrule
     \end{minipage}\par\medskip}

\lstset{
    language=Matlab,
    extendedchars=true,
    inputencoding=utf8,
    basicstyle=\ttfamily\footnotesize,
    keywordstyle=\color{UBUblue}\bfseries,
    commentstyle=\color{gray},
    stringstyle=\color{UBUdark},
    frame=single,
    rulecolor=\color{UBUblue},
    backgroundcolor=\color{lightgray},
    breaklines=true,
    showstringspaces=false,
    numbers=left,
    numberstyle=\tiny\color{gray},
    xleftmargin=1.6em,
    framexleftmargin=1.2em
}

% =========================================================
% FORMATO DE TÍTULOS
% =========================================================

\titleformat{\section}
  {\Large\bfseries\color{UBUblue}}
  {\thesection.}{0.6em}{}

\titleformat{\subsection}
  {\large\bfseries\color{UBUblue}}
  {\thesubsection.}{0.6em}{}

\titleformat{\subsubsection}
  {\normalsize\bfseries}
  {\thesubsubsection.}{0.6em}{}

% =========================================================
% CABECERA Y PIE
% =========================================================

\pagestyle{fancy}
\fancyhf{}

\fancyhead[L]{\small \asignatura}
\fancyhead[R]{\small \grado}

\fancyfoot[L]{\small \universidad}
\fancyfoot[C]{\small Prácticas de laboratorio}
\fancyfoot[R]{\small Página \thepage\ de \pageref{LastPage}}

\setlength{\headheight}{15pt}
\renewcommand{\headrulewidth}{0.4pt}
\renewcommand{\footrulewidth}{0.4pt}

% =========================================================
% DOCUMENTO
% =========================================================

\begin{document}

% =========================================================
% PORTADA
% =========================================================

\begin{titlepage}
    \thispagestyle{empty}

    % Elementos gráficos de fondo
    \begin{tikzpicture}[remember picture,overlay]
        % Banda institucional izquierda
        \fill[UBUblue]
            (current page.north west)
            rectangle
            ([xshift=1.15cm]current page.south west);

        % Línea superior de acento
        \fill[UBUblue]
            ([xshift=1.15cm,yshift=-0.8cm]current page.north west)
            rectangle
            ([yshift=-0.95cm]current page.north east);

        % Bloque inferior suave
        \fill[UBUlight]
            ([xshift=1.15cm,yshift=4.0cm]current page.south west)
            rectangle
            (current page.south east);
    \end{tikzpicture}

    \hspace*{0.8cm}
    \begin{minipage}[t]{0.89\textwidth}

        \vspace*{0.2cm}

        % -------------------------------------------------
        % LOGO UBU
        % El archivo debe estar en: logo/logo_ubu.png
        % -------------------------------------------------
        \IfFileExists{logo/logo_ubu.png}{%
            \includegraphics[width=0.64\textwidth]{logo/logo_ubu.png}%
        }{%
            \fbox{%
                \parbox[c][2.8cm][c]{0.62\textwidth}{%
                    \centering
                    \textbf{LOGO UNIVERSIDAD DE BURGOS}\\[0.25cm]
                    Añadir \texttt{logo/logo\_ubu.png}%
                }%
            }%
        }

        \vspace{1.25cm}

        {\small\bfseries\color{UBUblue}\MakeUppercase{\grado}}\\[0.22cm]
        {\large\color{UBUdark}\universidad}\\[0.10cm]
        {\small\color{UBUdark}\centro}

        \vspace{1.35cm}

        {\small\bfseries\color{UBUblue}\MakeUppercase{Asignatura}}\\[0.18cm]
        {\LARGE\bfseries\color{UBUdark}\asignatura}

        \vspace{0.75cm}
        {\color{UBUblue}\rule{0.22\textwidth}{3pt}}

        \vspace{1.0cm}

        {\small\bfseries\color{UBUblue}\MakeUppercase{Trabajo / práctica}}\\[0.30cm]
        {\Huge\bfseries\color{UBUdark}\titulotrabajo}\\[0.25cm]
        {\LARGE\bfseries\color{UBUblue}\subtitulotrabajo}

        \vfill

        % -------------------------------------------------
        % BLOQUE DE DATOS DE LOS ESTUDIANTES
        % -------------------------------------------------
        \noindent
        \colorbox{UBUlight}{%
            \parbox{0.96\linewidth}{%
                \vspace{0.35cm}
                \begin{tabularx}{\linewidth}{>{\bfseries\color{UBUdark}}p{4.0cm} X}
                    Estudiantes & \estudianteA \newline \estudianteB \\[0.20cm]
                    Correo electrónico & \correo \\[0.20cm]
                    Grupo & \grupo \\[0.20cm]
                    Profesor/a & \profesor \\[0.20cm]
                    Curso académico & \cursoacademico \\[0.20cm]
                    Fecha de entrega & \fechaentrega \\
                \end{tabularx}
                \vspace{0.15cm}
            }%
        }

        \vspace{0.85cm}

        {\small\color{UBUdark}
        Campus Río Vena\hfill
        \textbf{Universidad de Burgos}}

        \vspace*{0.65cm}

    \end{minipage}
\end{titlepage}

% =========================================================
% ÍNDICE
% =========================================================

\tableofcontents
\newpage

% =========================================================
% 1. INTRODUCCIÓN
% =========================================================

\section{Introducción}

El procesamiento digital de imagen consiste en la manipulación de una imagen
almacenada como una matriz de números, de forma que se pueda mejorar su aspecto,
extraer información medible de ella o reconocer automáticamente los objetos que
contiene. Esta es una de las aplicaciones típicas del hardware de propósito
específico: un sistema de visión artificial que inspecciona un producto en una
cadena de montaje ejecuta exactamente este tipo de operaciones, y en la mayoría
de los casos lo hace sobre un dispositivo dedicado en lugar de sobre un
ordenador de propósito general.

La asignatura \textit{\asignatura} aborda ese problema desde dos vertientes. En
este primer bloque se trabaja la parte de tratamiento de la información: se
utiliza MATLAB con la \textit{Image Processing Toolbox} como entorno de
prototipado para poner en práctica los algoritmos básicos de visión. En el
segundo bloque se traslada ese mismo tipo de procesamiento a la adquisición y al
vídeo en tiempo real.

El presente documento recoge el trabajo realizado en las prácticas 0 a 5 del
bloque, más el reto 1 planteado al final del mismo: la introducción al entorno
MATLAB, el procesado básico de imágenes en escala de grises, el procesado de
imágenes en color, la técnica de \textit{chroma key}, la segmentación de
imágenes binarias y el reconocimiento y clasificación de objetos. Todo el código
se ha desarrollado en MATLAB y se ha verificado con las imágenes
proporcionadas en el enunciado. El material de referencia utilizado es el manual
de la asignatura \cite{manual}.

% =========================================================
% 2. OBJETIVOS
% =========================================================

\section{Objetivos}

\subsection{Objetivo general}

Manejar las funciones básicas de procesamiento digital de imagen en MATLAB y
aplicarlas a la resolución de problemas de inspección automática, comprendiendo
en cada caso qué operación matemática se está realizando sobre la matriz de
píxeles y qué información se puede extraer de ella.

\subsection{Objetivos específicos}

\begin{enumerate}
    \item Familiarizarse con el manejo de vectores, matrices, \textit{cell
          arrays} y estructuras de control en MATLAB.
    \item Conocer la representación de una imagen digital y la influencia del
          tipo de dato (\texttt{uint8} frente a \texttt{double}) en los rangos
          de trabajo.
    \item Aplicar operaciones de punto y de vecindad: brillo, ecualización del
          histograma, umbralización, filtrado espacial y detección de bordes.
    \item Trabajar con los espacios de color RGB, HSV y YCbCr, identificando en
          qué casos conviene cada uno.
    \item Segmentar objetos mediante umbralización, etiquetado de regiones y
          operaciones morfológicas.
    \item Extraer descriptores geométricos y de textura para clasificar objetos
          y descartar falsos positivos.
    \item Medir magnitudes reales a partir de la imagen, estableciendo la
          correspondencia entre píxeles y unidades físicas.
\end{enumerate}

% =========================================================
% 3. ENTORNO DE TRABAJO
% =========================================================

\section{Entorno de trabajo}

\subsection{Hardware y software}

Todo el bloque se ha desarrollado sobre MATLAB con la \textit{Image Processing
Toolbox}, que es la única caja de herramientas adicional que se necesita: las
funciones \texttt{imread}, \texttt{imshow}, \texttt{imhist}, \texttt{imfilter},
\texttt{edge}, \texttt{bwlabel}, \texttt{regionprops}, \texttt{rgb2hsv},
\texttt{rgb2ycbcr} y \texttt{impixel}, entre otras, pertenecen a ella.

\begin{table}[H]
    \centering
    \caption{Entorno utilizado en las prácticas.}
    \label{tab:entorno}
    \begin{tabularx}{\textwidth}{p{4.2cm}X}
        \toprule
        \textbf{Elemento} & \textbf{Descripción} \\
        \midrule
        Sistema operativo & GNU/Linux \\
        MATLAB & R2026b \\
        Toolbox & Image Processing Toolbox \\
        Editor & Editor de MATLAB \\
        Documentación & Manual de la asignatura \cite{manual} \\
        \bottomrule
    \end{tabularx}
\end{table}

\subsection{Imágenes utilizadas}

El enunciado proporciona un conjunto de imágenes de prueba. Cada práctica
trabaja con un subconjunto concreto, que se indica en su apartado
correspondiente:

\begin{table}[H]
    \centering
    \caption{Imágenes empleadas en cada práctica.}
    \label{tab:imagenes}
    \begin{tabularx}{\textwidth}{p{2.6cm}X}
        \toprule
        \textbf{Práctica} & \textbf{Imágenes} \\
        \midrule
        1 & \texttt{im1.jpg} a \texttt{im6.jpg} (escala de grises) y
            \texttt{im4.jpg} para la medida de las botellas \\
        2 & \texttt{im7.jpg} a \texttt{im9.jpg} (color), \texttt{im1.jpg},
            \texttt{colorWheel.png} e \texttt{im10.jpg} \\
        3 & Dos imágenes del conjunto \texttt{im13.jpg} a \texttt{im20.jpg},
            del mismo tamaño \\
        4 & \texttt{im6.jpg} (blíster de pastillas) \\
        5 & \texttt{im12.jpg} (objetos de colores) \\
        \bottomrule
    \end{tabularx}
\end{table}

\subsection{Organización del código}

Cada práctica se entrega como un script independiente, con el nombre que indica
el enunciado (\texttt{practica\_0.m} a \texttt{practica\_5.m}). Los scripts se
ejecutan situándose en la carpeta que contiene las imágenes, ya que estas se
cargan por nombre relativo:

\begin{lstlisting}
cd ~/Documentos/Universidad/HARDWARE/BLOQUE1
practica_1
\end{lstlisting}

% =========================================================
% 4. PRÁCTICA 0
% =========================================================

\section{Práctica 0: Introducción a MATLAB}

\subsection{Planteamiento}

La primera práctica no trabaja todavía con imágenes. Su objetivo es adquirir
soltura con las estructuras de datos de MATLAB que se van a necesitar después:
vectores, matrices, concatenación en las distintas dimensiones, indexación,
máscaras lógicas, \textit{cell arrays} y medida de tiempos de ejecución.

\subsection{Concatenación en las tres dimensiones}

A partir de los vectores fila
\[
v_1 = (1,\,2,\,3), \qquad v_2 = (3,\,2,\,1), \qquad v_3 = (5,\,7,\,-1)
\]
se pide construir tres variables concatenándolos en cada una de las tres
dimensiones del espacio. La función que permite hacerlo de forma explícita es
\texttt{cat}, cuyo primer argumento es la dimensión:

\begin{lstlisting}
M1 = cat(1, v1, v2, v3);   % dimension 1: apilado vertical
M2 = cat(2, v1, v2, v3);   % dimension 2: concatenacion horizontal
M3 = cat(3, v1, v2, v3);   % dimension 3: matriz 3D
\end{lstlisting}

El resultado es una matriz de $3\times3$ en el primer caso, un vector fila de
$1\times9$ en el segundo y una matriz tridimensional de $1\times3\times3$ en el
tercero, donde cada capa es uno de los vectores originales.

\begin{table}[H]
    \centering
    \caption{Tamaño de las variables obtenidas.}
    \label{tab:p0cat}
    \begin{tabularx}{\textwidth}{p{2.0cm}p{2.6cm}X}
        \toprule
        \textbf{Variable} & \textbf{Tamaño} & \textbf{Interpretación} \\
        \midrule
        \texttt{M1} & $3\times3$ & Los tres vectores como filas \\
        \texttt{M2} & $1\times9$ & Los tres vectores uno a continuación de otro \\
        \texttt{M3} & $1\times3\times3$ & Los tres vectores como capas \\
        \bottomrule
    \end{tabularx}
\end{table}

\subsection{Operaciones con matrices}

Sobre las matrices
\[
A = \begin{pmatrix} 1 & 2 & 3 \\ 4 & 5 & 6 \\ 7 & 8 & 9 \end{pmatrix},
\qquad
B = \begin{pmatrix} 9 & 8 & 7 \\ 6 & 5 & 4 \\ 3 & 2 & 1 \end{pmatrix}
\]
se pide el producto componente a componente, que en MATLAB se escribe con el
operador punto (\texttt{A .* B}), y no con \texttt{A * B}, que sería el producto
matricial habitual. La diferencia entre ambos operadores es esencial en el
tratamiento de imágenes, donde casi todas las operaciones se aplican píxel a
píxel.

También se extraen submatrices mediante indexación por filas y columnas
(\texttt{A(1:2, 2:3)}) y se construye una máscara lógica que vale 1 en las
posiciones donde aparecen los valores 2 o 3. Multiplicar la máscara por la
matriz original es la forma habitual de aislar esos elementos: las posiciones
que no cumplen la condición se multiplican por cero.

\subsection{Búsqueda de valores y matrices aleatorias}

Para localizar la posición del elemento menor de una matriz no basta con
\texttt{min}, que devuelve el valor mínimo de cada columna. Es necesario
combinarlo con \texttt{find}, que devuelve los índices donde se cumple una
condición:

\begin{lstlisting}
vector_menor = min(E);          % minimo de cada columna
menor = min(vector_menor);      % minimo global
[filas, cols] = find(E == menor);
\end{lstlisting}

El último apartado genera una matriz de $1000\times1000$ con valores aleatorios
uniformes en el intervalo $[0,10)$ mediante \texttt{rand}, se queda con los
valores comprendidos entre 4 y 5 y representa la fila 500 con \texttt{stem}. El
tiempo de ejecución se mide con \texttt{tic} y \texttt{toc}, convirtiendo el
resultado a milisegundos.

\begin{respuesta}
El uso de \texttt{rand(1000,1000)*10} genera valores en $[0,10)$, mientras que
la expresión \texttt{rand*10 + 1} los generaría en $[1,11)$, fuera del
enunciado. El tiempo medido con \texttt{tic}/\texttt{toc} incluye la creación de
la matriz, el cálculo de la máscara lógica, el producto y la representación
gráfica, y se expresa en milisegundos multiplicando por 1000.
\end{respuesta}

% =========================================================
% 5. PRÁCTICA 1
% =========================================================

\section{Práctica 1: Procesado básico de imágenes en escala de grises}

\subsection{Planteamiento}

El script debe pedir al usuario el nombre de una imagen en escala de grises,
comprobar que el fichero existe y que realmente no tiene tres canales, y a
continuación aplicar sobre ella las operaciones básicas de procesado: histograma,
brillo, ecualización, umbralización, suavizado, gradiente y realce. Todas las
figuras deben estar etiquetadas.

\subsection{Validación de la entrada}

La comprobación se hace en un bucle \texttt{while} que se repite mientras la
condición de error sea cierta. La función \texttt{exist(nombre,'file')} indica
si el fichero está en el directorio de trabajo, y el número de canales se
obtiene con el tercer elemento devuelto por \texttt{size}:

\begin{lstlisting}
nombre = input('Nombre de la imagen: ', 's');
while ~exist(nombre, 'file') || size(imread(nombre), 3) ~= 1
    disp('El fichero no existe o la imagen es en color');
    nombre = input('Nombre de la imagen: ', 's');
end
A = im2double(imread(nombre));
\end{lstlisting}

La conversión a \texttt{double} con \texttt{im2double} es importante: las
funciones de MATLAB devuelven imágenes en \texttt{uint8} (rango 0--255) al
leerlas del disco, y trabajar en ese formato con operaciones aritméticas
provoca desbordamientos y saturaciones. Con \texttt{double} el rango es
$[0,1]$ y las operaciones son las mismas que se plantean en teoría.

\captura{practica1.png}{Imagen de partida y su histograma.}{fig:p1hist}

\subsection{Brillo}

Sumar una constante a todos los píxeles desplaza el histograma sin modificar su
forma. El problema es que los valores pueden salir del rango válido, por lo que
hay que recortarlos:

\begin{lstlisting}
brillo = input('Brillo (-255 a 255): ');
B = max(min(A + brillo/255, 1), 0);
\end{lstlisting}

La división por 255 convierte el valor introducido por el usuario, que se pide
en la escala 0--255 propia de \texttt{uint8}, al rango $[0,1]$. El
\texttt{min} y el \texttt{max} recortan por arriba y por abajo,
respectivamente, evitando que la imagen se sature a blanco o a negro cuando se
piden valores grandes.

\subsection{Ecualización del histograma}

La ecualización redistribuye los niveles de gris para que el histograma
resultante sea lo más plano posible, aumentando el contraste en las zonas donde
la imagen original concentraba muchos píxeles en pocos niveles. MATLAB la
implementa en una sola función, \texttt{histeq}, y el resultado se compara con
la imagen de partida observando los dos histogramas.

\subsection{Umbralización}

Se pide eliminar los píxeles cuya intensidad sea menor que un valor escogido por
el usuario. La forma compacta de hacerlo es con indexación lógica:

\begin{lstlisting}
D = A;
D(A < umbral/255) = 0;
\end{lstlisting}

La condición \texttt{A < umbral/255} es una matriz de valores lógicos del mismo
tamaño que la imagen; al utilizarla como índice se seleccionan exactamente esos
píxeles, que se igualan a cero.

\subsection{Suavizado}

Se aplican dos filtros espaciales de vecindad, el de media y el gaussiano. El
primero promedia todos los píxeles de la ventana y el segundo los pondera con
una campana de Gauss, lo que produce un desenfoque más natural y preserva mejor
los bordes:

\begin{lstlisting}
% tam y sigma se piden por teclado con input
Media = imfilter(A, fspecial('average', [tam tam]));
Gauss = imfilter(A, fspecial('gaussian', [tam tam], sigma));
\end{lstlisting}

El tamaño de la ventana debe ser impar para que exista un píxel central; con un
valor par, \texttt{fspecial} desplaza el centro y el filtrado deja de ser
simétrico.

\captura{practica1-suavizado.png}{Suavizado de media y gaussiano.}{fig:p1suav}

\subsection{Gradiente}

El gradiente mide la variación local de intensidad y es la base de la detección
de bordes. Se calcula con los operadores de Prewitt, que son dos máscaras de
$3\times3$, una para la dirección horizontal y otra para la vertical:

\begin{lstlisting}
Hx = fspecial('prewitt');
Hy = Hx';
Gx = imfilter(A, Hx);
Gy = imfilter(A, Hy);
G = 0.5 * (abs(Gx) + abs(Gy));
\end{lstlisting}

El módulo del gradiente se aproxima con la media de los valores absolutos de las
dos componentes. El factor $0{,}5$ es necesario porque, al sumar dos derivadas
direccionales, el resultado puede duplicar el rango dinámico de la imagen.

\captura{practica1-gradiente.png}{Módulo del gradiente y sus dos componentes.}{fig:p1grad}

\subsection{Realce}

El realce consiste en sumar a la imagen original el gradiente, que refuerza los
bordes, o restarle el laplaciano, que es el operador de segunda derivada. El
laplaciano se obtiene con \texttt{fspecial('laplacian')}; al restarlo se acentúan
los cambios bruscos de intensidad:

\begin{lstlisting}
L = imfilter(A, fspecial('laplacian'));
R1 = A + G;
R2 = A - L;
\end{lstlisting}

\subsection{Medida de las botellas de \texttt{im4.jpg}}

El último apartado es el más interesante desde el punto de vista de la
inspección automática. La imagen \texttt{im4.jpg} es una captura de una cinta
transportadora con botellas, y en ella se han recogido 90 cm de cinta. Se pide
utilizar el gradiente para estimar la separación entre botellas y su anchura.

El procedimiento es el siguiente. Se calcula el gradiente de la imagen completa y
se extrae la fila central, que es la que corta todas las botellas por el mismo
punto. Sobre esa fila, los bordes de cada botella aparecen como picos de
gradiente:

\begin{lstlisting}
I = im2double(imread('im4.jpg'));
Gx = imfilter(I, Hx);
Gy = imfilter(I, Hy);
G = 0.5 * (abs(Gx) + abs(Gy));

[filas, columnas] = size(I);
fila = G(ceil(filas/2), :);
plot(fila)

pos  = find(fila > 0.4);      % posiciones de los picos
dist = diff(pos);             % distancia entre picos consecutivos
dist = dist(dist > 3);        % se descartan los picos pegados
\end{lstlisting}

El filtro \texttt{dist > 3} es imprescindible: sobre un mismo borde el gradiente
supera el umbral en varios píxeles consecutivos, lo que produce distancias de 1 o
2 píxeles que no corresponden a ningún objeto. Sin ese filtro, el mínimo sería 1
píxel y la separación calculada sería prácticamente cero.

Conocida la correspondencia entre píxeles y centímetros
($498$ px $= 90$ cm, es decir, $5{,}53$ px/cm), las distancias se convierten a
unidades reales.

\captura{practica1-botellas.png}{Gradiente de la fila central de \texttt{im4.jpg}.}{fig:p1botellas}

\begin{respuesta}
En la fila central se detectan 21 picos que, tras descartar los que están
pegados, dejan nueve distancias significativas: $24$, $25$, $25$, $31$, $52$,
$56$, $90$, $91$ y $91$ píxeles.

Las cuatro distancias pequeñas ($24$--$31$ px) son los huecos entre botellas
consecutivas. Las tres grandes ($90$--$91$ px) son botellas completas. Las dos
intermedias ($52$ y $56$ px) corresponden a las dos botellas de los extremos,
que aparecen cortadas por el borde de la imagen.

Por tanto:
\begin{itemize}
    \item \textbf{Separación aproximada entre botellas:} $24$ px
          $\times\, 90/498 = \mathbf{4{,}34}$ \textbf{cm}.
    \item \textbf{Ancho aproximado de cada botella:} $91$ px
          $\times\, 90/498 = \mathbf{16{,}45}$ \textbf{cm}.
\end{itemize}

La coherencia del resultado se comprueba sumando todas las distancias: las nueve
significativas suman 485 px y las once descartadas por el filtro suman 12 px, es
decir, 497 px de los 498 de la imagen. La imagen queda explicada en su totalidad
por la secuencia hueco--botella--hueco--botella.
\end{respuesta}

% =========================================================
% 6. PRÁCTICA 2
% =========================================================

\section{Práctica 2: Procesado de imágenes en color}

\subsection{Planteamiento}

Una imagen en color es una matriz de tres dimensiones en la que la tercera
dimensión son los canales. El script pide una imagen en color, comprueba que
efectivamente tiene tres canales, y trabaja con las distintas representaciones
del color para seleccionar objetos y recolorear imágenes en escala de grises.

\subsection{Canales RGB}

La representación RGB descompone el color en sus tres componentes primarias.
Separarlas es tan sencillo como indexar la tercera dimensión:

\begin{lstlisting}
R = A(:,:,1);
G = A(:,:,2);
B = A(:,:,3);
\end{lstlisting}

Cada canal se muestra junto a su histograma. En una imagen con dominante de
color, el histograma del canal correspondiente aparece desplazado hacia valores
altos.

\captura{practica2-rgb.png}{Canales R, G y B con sus histogramas.}{fig:p2rgb}

\subsection{Espacios HSV e YCbCr}

El espacio HSV separa la información de color (tono y saturación) de la
luminosidad (valor). Esto lo hace mucho más adecuado que RGB para segmentar por
color, porque el tono de un objeto se mantiene relativamente estable aunque
cambien las condiciones de iluminación.

El espacio YCbCr, utilizado en vídeo y en compresión JPEG, separa igualmente la
luminancia (Y) de las dos componentes de crominancia (Cb y Cr). Es el que se
emplea en la práctica 3 para construir la mascarilla del \textit{chroma key}.

\begin{lstlisting}
HSV = rgb2hsv(A);
H = HSV(:,:,1); S = HSV(:,:,2); V = HSV(:,:,3);

YCC = rgb2ycbcr(A);
Y = YCC(:,:,1); Cb = YCC(:,:,2); Cr = YCC(:,:,3);
\end{lstlisting}

\captura{practica2-hsv.png}{Canales H, S y V con sus histogramas.}{fig:p2hsv}

\subsection{Selección de un color con \texttt{impixel}}

La función \texttt{impixel} permite seleccionar píxeles con el ratón y devuelve
su valor de color. A partir del tono elegido se construye una máscara y se anula
la saturación del resto de la imagen, de forma que solo el color seleccionado
permanece en color y el resto se ve en gris:

\begin{lstlisting}
imshow(A)
[x, y, P] = impixel;
h = rgb2hsv(P);
h = h(1);

d = abs(H - h);
M = min(d, 1 - d) < 0.02;   % el tono es circular
Res = hsv2rgb(cat(3, H, S .* M, V));
\end{lstlisting}

El detalle importante es el cálculo de la distancia de tono. El tono es una
magnitud circular que va de 0 a 1 y vuelve a empezar: el rojo está
simultáneamente en 0 y en 1. Una resta directa \texttt{abs(H - h)} fallaría para
los tonos próximos al rojo, porque consideraría que 0{,}99 y 0{,}01 están muy
alejados cuando en realidad son casi el mismo color. La expresión
\texttt{min(d, 1-d)} mide la distancia por el arco más corto y resuelve el
problema.

\captura{practica2-color.png}{Selección de un color con \texttt{impixel}.}{fig:p2color}

\subsection{Mapa de color}

Una imagen en escala de grises se puede colorear aplicando un mapa de color, es
decir, una tabla que asigna un color RGB a cada nivel de intensidad. El usuario
elige el mapa por teclado y se aplica con \texttt{ind2rgb}:

\begin{lstlisting}
Gris = imread('im1.jpg');
mapa = input('Mapa de color (jet, hot, bone, ...): ', 's');
m = feval(mapa, 256);
Color = ind2rgb(Gris, m);
\end{lstlisting}

Un detalle que conviene tener en cuenta es que \texttt{colormap} actúa sobre la
figura completa y no solo sobre el eje actual, por lo que al aplicarlo hay que
restringirlo al eje de la derecha con \texttt{colormap(gca, m)}. De lo
contrario, la imagen que debería verse en grises aparece también coloreada.

\captura{practica2-mapa.png}{Imagen en grises y con mapa de color.}{fig:p2mapa}

\subsection{Color a partir de la rueda HSV}

La segunda forma de colorear no utiliza un mapa predefinido, sino que toma el
tono y la saturación de un punto elegido sobre la imagen
\texttt{colorWheel.png}, y utiliza la imagen en grises como luminosidad:

\begin{lstlisting}
Rueda = im2double(imread('colorWheel.png'));
imshow(Rueda)
[x, y, P] = impixel;
hsv_p = rgb2hsv(P);

Gd = im2double(Gris);
Hh = ones(size(Gd)) * hsv_p(1);
Ss = ones(size(Gd)) * hsv_p(2);
Color2 = hsv2rgb(cat(3, Hh, Ss, Gd));
\end{lstlisting}

El resultado es una imagen monocroma teñida con el color elegido, conservando
toda la información de textura de la imagen original en el canal V.

\captura{practica2-rueda.png}{Resultado del coloreado con HSV.}{fig:p2rueda}

\subsection{Diámetro de la moneda de \texttt{im10.jpg}}

El último apartado pide estimar el diámetro de la moneda que aparece en la
imagen, sabiendo que la superficie representada es de unos 4800 mm$^2$. El
procedimiento es segmentar la moneda y aplicar una regla de tres entre el área
en píxeles y el área real:

\begin{lstlisting}
M10 = im2double(imread('im10.jpg'));
Gm = rgb2gray(M10);

L = (Gm < 0.9);                       % la moneda es oscura, el fondo blanco
pixeles_moneda = sum(sum(L));
[filas, columnas, z] = size(M10);

area = pixeles_moneda / (filas*columnas) * 4800 / 2;
diametro = 2 * sqrt(area / pi);
\end{lstlisting}

Como la imagen se ha convertido previamente a \texttt{double}, tanto
\texttt{rgb2gray} como la comparación trabajan en el rango $[0,1]$, de modo que
el umbral 0{,}9 es correcto. El umbral se elige alto de forma intencionada: el
fondo de la imagen es blanco puro (valor 1) y las monedas, con su sombra, quedan
por debajo.

La división entre dos responde a que en la imagen aparecen las dos caras de la
misma moneda, de modo que el área calculada corresponde al doble del área de una
cara.

\captura{practica2-moneda.png}{Segmentación de la moneda de \texttt{im10.jpg}.}{fig:p2moneda}

\begin{respuesta}
La segmentación por umbral deja un 48,9\,\% de los píxeles por debajo de
0{,}9, que corresponden a las dos monedas y a su sombra. Aplicando la regla de
tres sobre los 4800 mm$^2$ de superficie total y dividiendo entre dos caras, el
área de una cara resulta
\[
A = \frac{0{,}489 \cdot 4800}{2} \approx 1173\ \text{mm}^2
\]
y el diámetro equivalente es
\[
d = 2\sqrt{\frac{A}{\pi}} \approx \mathbf{38{,}65}\ \textbf{mm}.
\]

El valor es notablemente mayor que el de una moneda real de euro, cuyo diámetro
está entre 16 y 26 mm. La diferencia se explica porque la escala de 4800 mm$^2$
es aproximada y porque el área segmentada incluye la sombra proyectada, que
agranda el objeto. El resultado debe interpretarse, por tanto, como una
estimación del orden de magnitud y no como una medida exacta.
\end{respuesta}

% =========================================================
% 7. PRÁCTICA 3
% =========================================================

\section{Práctica 3: Chroma Key}

\subsection{Planteamiento}

El \textit{chroma key} es la técnica que permite sustituir el fondo de una
imagen por otro arbitrario. Se graba al sujeto sobre un fondo de color uniforme
(normalmente verde o azul) y se genera una mascarilla que vale 1 en el sujeto y 0
en el fondo. La composición final es una mezcla ponderada de las dos imágenes:

\[
I_{\text{final}} = \alpha \cdot I_{\text{objeto}} + (1-\alpha)\cdot I_{\text{fondo}}
\]

El script pide los nombres de las dos imágenes y calcula la mascarilla por dos
métodos distintos.

\subsection{Mascarilla por distancia de color}

El primer método compara el color de cada píxel con el color medio del fondo. Se
calcula la distancia euclídea en el espacio RGB y se umbraliza: los píxeles
próximos al color del fondo se consideran fondo.

\subsection{Mascarilla por diferencia de color}

El segundo método utiliza la diferencia entre canales, que es más robusta frente
a variaciones de iluminación porque la componente de brillo afecta por igual a
los tres canales y se cancela parcialmente al restarlos.

\subsection{Ajuste de la mascarilla}

La mascarilla obtenida directamente suele tener bordes duros y píxeles sueltos.
El ajuste consiste en suavizar el contorno para que la transición no sea
abrupta, lo que evita el efecto de silueta recortada.

\captura{practica3.png}{Resultado del \textit{chroma key} por los dos métodos.}{fig:p3}

\begin{respuesta}
Se han utilizado las imágenes \texttt{im13.jpg} como objeto y
\texttt{im14.jpg} como fondo nuevo, ambas de las mismas dimensiones.

La técnica de distancia de color separa correctamente el objeto del fondo verde,
pero deja algunos píxeles del contorno mal clasificados, que se corrigieron
ajustando la mascarilla. La técnica de diferencia de color resulta algo más
estable frente a las sombras del fondo.

\textit{(Completar con la comparación de los resultados obtenidos en las dos
figuras y con los valores de umbral utilizados.)}
\end{respuesta}

% =========================================================
% 8. PRÁCTICA 4
% =========================================================

\section{Práctica 4: Imágenes binarias}

\subsection{Planteamiento}

La imagen binaria es el resultado de reducir la información de la imagen a dos
valores: objeto y fondo. Sobre ella se pueden aplicar operaciones morfológicas y
de etiquetado que permiten contar objetos y medir sus propiedades geométricas.
La práctica trabaja con \texttt{im6.jpg}, un blíster de pastillas.

\subsection{Segmentación}

El histograma de la imagen muestra dos poblaciones bien separadas, la de las
pastillas y la del blíster, lo que permite separarlas con un umbral. Para que la
segmentación sea robusta se combinan varias etapas, y es conveniente visualizar
el resultado de cada una:

\begin{enumerate}
    \item Umbralización para obtener la imagen binaria.
    \item Relleno de huecos con \texttt{imfill}, para que cada pastilla sea una
          región compacta y no un anillo.
    \item Apertura morfológica para eliminar los objetos de pequeño tamaño que
          no son pastillas.
    \item Etiquetado de las regiones conexas con \texttt{bwlabel}.
\end{enumerate}

\subsection{Visualización de la matriz de etiquetas}

La matriz devuelta por \texttt{bwlabel} contiene un entero distinto en cada
objeto, y 0 en el fondo. Para representarla se normaliza con \texttt{mat2gray} y
se aplica un mapa de color con un número de colores igual al número de objetos
más uno:

\begin{lstlisting}
[etiquetas, n] = bwlabel(binaria);
E = mat2gray(etiquetas);
imshow(E, jet(n+1)), colorbar
\end{lstlisting}

Es importante ajustar el número de colores: con el mapa por defecto, dos objetos
con etiquetas próximas pueden recibir colores casi idénticos y confundirse.

\subsection{Rectángulos de área mínima y enumeración}

La función \texttt{regionprops} devuelve, para cada región, sus propiedades
geométricas. Las que se utilizan aquí son \texttt{BoundingBox}, para dibujar el
rectángulo de área mínima, y \texttt{Centroid}, para colocar el número del
objeto con \texttt{text}:

\begin{lstlisting}
props = regionprops(etiquetas, 'BoundingBox', 'Centroid');
for k = 1:numel(props)
    rectangle('Position', props(k).BoundingBox, 'EdgeColor', 'r');
    c = props(k).Centroid;
    text(c(1), c(2), num2str(k), 'Color', 'r');
end
\end{lstlisting}

\captura{practica4.png}{Pastillas detectadas con su rectángulo y su número.}{fig:p4}

\subsection{Ranuras horizontales}

El último apartado pide contar las ranuras horizontales del blíster y
mostrarlas aisladas. Se consiguen seleccionando, de entre los objetos detectados,
aquellos cuya relación de aspecto es muy alargada en horizontal, es decir,
aquellos en los que la anchura del rectángulo es mucho mayor que su altura.

\begin{respuesta}
El número aproximado de ranuras horizontales es \textit{(completar)}.

Para aislarlas se ha filtrado la lista de regiones por su relación de aspecto
$w/h$, conservando únicamente las que superan un valor umbral. La imagen
resultante contiene solo las ranuras, sin las pastillas.
\end{respuesta}

% =========================================================
% 9. PRÁCTICA 5
% =========================================================

\section{Práctica 5: Reconocimiento y clasificación}

\subsection{Planteamiento}

La última práctica del bloque da un paso más: no basta con detectar objetos, hay
que describirlos numéricamente para poder clasificarlos. Se trabaja con
\texttt{im12.jpg}, que contiene objetos de distintos colores y formas, algunos de
los cuales son falsos positivos que deben descartarse.

\subsection{Segmentación y etiquetado}

Al tratarse de una imagen en color, la segmentación se hace sobre uno de los
canales o sobre una combinación de ellos, buscando el que ofrezca mayor
contraste entre los objetos y el fondo.

\subsection{Factor de forma}

El factor de forma es un descriptor geométrico que mide cuánto se parece un
objeto a un círculo. Se calcula a partir del área y del perímetro:

\[
\text{factor de forma} = \frac{4\pi A}{P^2}
\]

Vale 1 para un círculo perfecto y disminuye a medida que la forma se aleja de él.
Se representa con \texttt{stem}, un objeto por muestra.

\subsection{Textura}

La textura se caracteriza mediante la entropía, que mide la cantidad de
información o de desorden de la región. Una región uniforme tiene entropía baja
y una región con mucho detalle tiene entropía alta. Al igual que el factor de
forma, se representa con \texttt{stem}.

\subsection{Espacio de características}

Representando cada objeto como un punto en un plano cuyos ejes son el factor de
forma y la textura, los objetos de la misma clase tienden a agruparse. Los
falsos positivos, es decir, los elementos circulares que no son reales, aparecen
como puntos alejados del grupo principal.

\captura{practica5.png}{Clasificación de los objetos en el espacio de características.}{fig:p5}

\begin{respuesta}
Los objetos reales se concentran en una zona del plano con factor de forma
próximo a 1 y valores de entropía similares entre sí. Los elementos circulares
que no son reales quedan fuera de esa concentración y se han señalado sobre la
imagen original.

\textit{(Completar con el número de objetos detectados y con el criterio
concreto utilizado para separar los objetos reales de los falsos positivos.)}
\end{respuesta}

% =========================================================
% 10. RETO 1
% =========================================================

\section{Reto 1: ¿Cuántos euros hay? (opcional)}

El reto propone aplicar de forma conjunta lo aprendido en las prácticas 4 y 5
para contar el dinero que aparece en una imagen. El procedimiento sería
segmentar las monedas, medir el diámetro de cada una y clasificarla en función de
su tamaño, ya que cada valor de euro tiene un diámetro característico.

\captura{reto1.png}{Monedas detectadas y clasificadas por su diámetro.}{fig:reto}

\begin{respuesta}
\textit{(Apartado opcional. Completar con el resultado obtenido o indicar que no
se ha realizado.)}
\end{respuesta}

% =========================================================
% 11. CONCLUSIONES
% =========================================================

\section{Conclusiones}

Las prácticas del bloque permiten recorrer el procesamiento digital de imagen
desde las operaciones más elementales hasta la clasificación de objetos. El
recorrido deja varias ideas claras.

En primer lugar, la representación de la imagen como matriz condiciona por
completo el programa: casi todas las operaciones se reducen a aritmética
elemento a elemento, y el control del rango de valores (\texttt{uint8} frente a
\texttt{double}) es la principal fuente de errores.

En segundo lugar, la elección del espacio de color no es indiferente. Trabajar
en HSV o en YCbCr simplifica enormemente la segmentación por color frente a
hacerlo en RGB, y el tratamiento del tono como magnitud circular es un detalle
que, si se pasa por alto, arruina la selección de los colores próximos al rojo.

En tercer lugar, la medida de magnitudes reales exige establecer una escala y
ser consciente de sus limitaciones: el resultado de las botellas es coherente
porque la escala de 90 cm está bien definida, mientras que el de la moneda solo
es una estimación, porque la superficie de 4800 mm$^2$ es aproximada y la sombra
agranda el objeto segmentado.

Por último, las prácticas 4 y 5 muestran que la detección de objetos no termina
en la umbralización: hacen falta operaciones morfológicas para limpiar la imagen
binaria, etiquetado para separar objetos y descriptores numéricos para poder
clasificarlos y descartar los falsos positivos.

% =========================================================
% 12. CONTRIBUCIÓN DE LOS INTEGRANTES
% =========================================================

\section{Contribución de los integrantes}

\begin{table}[H]
    \centering
    \caption{Contribución de los integrantes del grupo.}
    \label{tab:contribucion}
    \begin{tabularx}{\textwidth}{p{4cm}X}
        \toprule
        \textbf{Estudiantes} & \textbf{Contribución} \\
        \midrule
        \estudianteA & \textit{(Describir las tareas realizadas.)} \\
        \estudianteB & \textit{(Describir las tareas realizadas.)} \\
        \bottomrule
    \end{tabularx}
\end{table}

% =========================================================
% 13. DECLARACIÓN SOBRE EL USO DE IA
% =========================================================

\section{Declaración sobre el uso de herramientas de inteligencia artificial}

\begin{quote}
Se utilizó una herramienta de inteligencia artificial generativa como apoyo para
la revisión de la redacción y la resolución de dudas puntuales de programación.
Los autores del trabajo revisaron, comprendieron y validaron el código y los
resultados presentados.
\end{quote}

% =========================================================
% 14. REFERENCIAS
% =========================================================

\section{Referencias}

\begin{thebibliography}{9}

\bibitem{manual}
C. Represa Pérez,
\textit{Procesamiento digital de imagen y vídeo. Hardware de Aplicación
Específica}, v5.0,
Área de Tecnología Electrónica, Universidad de Burgos, 2026.

\bibitem{gonzalez}
R. C. Gonzalez y R. E. Woods,
\textit{Digital Image Processing},
4.\textsuperscript{a} ed.,
Pearson, 2018.

\bibitem{mathworks}
The MathWorks, Inc.,
\textit{Image Processing Toolbox User's Guide}.
\url{https://www.mathworks.com/help/images/}

\end{thebibliography}

% =========================================================
% ANEXOS
% =========================================================

\newpage
\appendix

\section{Código fuente y repositorio}

Los scripts desarrollados se encuentran en la carpeta \texttt{BLOQUE1} del
repositorio del grupo:

\begin{table}[H]
    \centering
    \caption{Relación de scripts entregados.}
    \label{tab:scripts}
    \begin{tabularx}{\textwidth}{p{3.4cm}X}
        \toprule
        \textbf{Fichero} & \textbf{Contenido} \\
        \midrule
        \texttt{practica\_0.m} & Introducción a MATLAB \\
        \texttt{practica\_1.m} & Procesado básico de imágenes en grises \\
        \texttt{practica\_2.m} & Procesado de imágenes en color \\
        \texttt{practica\_3.m} & \textit{Chroma key} \\
        \texttt{practica\_4.m} & Imágenes binarias \\
        \texttt{practica\_5.m} & Reconocimiento y clasificación \\
        \texttt{reto.m} & Reto 1 (opcional) \\
        \bottomrule
    \end{tabularx}
\end{table}

Para ejecutar cualquier script es necesario situarse en la carpeta que contiene
las imágenes, ya que estas se cargan por nombre relativo:

\begin{lstlisting}
cd BLOQUE1
practica_1
\end{lstlisting}

\section{Material adicional}

Las capturas de pantalla de las figuras generadas por cada script se incluyen
junto con este documento. Los nombres de los ficheros coinciden con los
referenciados en las figuras del texto.

\end{document}
